\chapter{Testing}
In questa sezione vedremo quali test ho eseguito per verificare la correttezza degli algoritmi.
Per implementare i test ho deciso di utilizzare la libreria unittest di python con
eventualmente dei mock presi sempre dalla libreria unittest.

\section{Cifratura}
Per quanto riguarda la cifratura ho deciso di eseguire unit tests su tutti gli algoritmi
implementati. In particolare ho eseguito test di cifratura e decifratura per DES e RSA. Inoltre ho
testato la verifica della firma per RSA e il corretto scambio di chiavi per DH.

\section{Altro}

\subsection{Handshake}
Ho eseguito unittest anche per quanto riguarda Client e Server Handshake per verificare il corretto
scambio delle chiavi e la corretta verifica della firma.

\subsection{Connection Handler}
Ho deciso di testare il ConnectionHandler per verificare il corretto invio e ricezione di messaggi e
anche il corretto scambio di chiavi tramite handshake.

\subsection{Message router}
Ovviamente ho testato anche il router per verificare il corretto routing.

\subsection{Test di accesso non autorizzato}
Fortunatamente ho pensato all'idea di manomettere l'invio dei messaggi richiedendo messaggi di
altri utenti o inviando messaggi facendo finta di essere altri senza aver fatto il login. Una volta
testato mi sono accorto che il mio codice aveva un bug perche' non manteneva correttamente le
sessioni e non veniva controllato che il login doveva essere fatto dall'utente di cui si
richiedevano i messaggi. Quindi e' presente anche un test per controllare gli accessi non
autorizzati.
